\section{Conclusion}
\label{sec:conclusion}
We asked why splitting a function-node, which leaves the objective unchanged,
changes what Min-sum does.  On a single split binary function-node the answer
is exact: the two copies exchange messages through a short loop, the loop
absorbs the unary influence of the rest of the graph, and the variables commit
to one assignment within a few iterations (Section~\ref{sec:setting}).
Without damping this commitment is too fast.  Under the selection recursion,
which we verified on every audited instance, the run either converges or
alternates between exactly two coherent solutions, and on bipartite graphs a
rephasing turns the alternation into a fixed point
(Section~\ref{sec:tworoutes}).
Damping above an instance-dependent threshold removes the alternation, and
no problem-independent threshold exists (Section~\ref{sec:whydamping}).
Selecting between the two branches with standard local search recovers most
of the gap to the damped algorithm, which confirms the two-solution reading
but does not replace damping.

Empirically, splitting plus damping is what makes damped Min-sum settle.  On
the random families it reaches the same final cost as damped Min-sum alone
but settles about ten times earlier, and on the structured benchmarks it also
lowers the final cost.  Delaying the split improves on immediate splitting
on every family; with held-out selection of the delay its gain is established
on random dense and scale-free graphs.  The split ratio is a second knob: a
fixed $0.95/0.05$ split closes the gap to the learned solver DABP on the
random families without any learning but hurts on the structured ones, and a
temporary $0.95/0.05$ pulse takes most of the former gain without the latter
loss.  DABP remains ahead only on random dense graphs, by $0.1\%$, at the
price of centralized execution.

Two questions remain open.  The selection recursion is verified, not proved,
and a proof would turn the two-solution theorem into an unconditional
statement about Min-sum on split graphs.  The split ratio and the split time
are the only two parameters of the method; both were fixed by hand here, and
choosing them from the messages themselves is the natural next step.
